\documentclass[12pt]{exam}
\usepackage{amsmath,amsthm,amsbsy,amsfonts}
 
\boxedpoints
\pointsinmargin
\addpoints
 
\newcommand{\Z}{\mathbb Z}
\newcommand{\E}{\mathbb E}
\newcommand{\R}{\mathbb R}
\newcommand{\Q}{\mathbb Q}
 
\pagestyle{headandfoot}
%\extraheadheight{.4in}
\extrafootheight{-1in}
\extrawidth{.7in}
\firstpageheader{}{\bf\Large Brief Review}{}
 
\begin{document}
%Complete the following review questions before the next class. % Compare 
%your work with the answers at 
%\textit{http://www.math.ttu.edu/$\sim$williams/1352/review-answers.pdf}.

\begin{questions}

\question Consider a particle moving with velocity $v(t) = \sqrt{t}+3$ ft/sec from $t=0$ to $t=3$ seconds.
\begin{parts}
\part Make a very rough guess about the distance traveled by the particle.
\part Use a Riemann sum with 4 subdivisions and left endpoints to estimate this distance.
\part Repeat using 6 subdivisions.  What do you notice?
\part Use a definite integral to find the exact distance traveled by the particle.  Compare with your
estimates.
\end{parts}

\question Consider the area under the graph of $y=\sqrt{x}+3$ between $x=0$ and $x=2$.
\begin{parts}
\part Use a Riemann sum with 4 subdivisions and left endpoints to estimate this area.  Draw the region
and the 4 rectangles.
\part Use a definite integral to find the exact area.
\part Explain the connection between Riemann sums, definite integrals, distance, and area.
\end{parts}


\question Find $\displaystyle \frac {dy}{dx}$ if
\begin{parts}
\part $\displaystyle y=x^7-5x^{\frac 3 2} + \frac{4}{x} -3$
\vspace{0.3in}

\part $\displaystyle y=\sqrt{\cos\sqrt x}$
\vspace{0.3in}

\part $\displaystyle y=e^{4x}\tan{x^2}$
\vspace{0.3in}

\part $\displaystyle y=\frac{3x^2\csc x}{\sin^{-1}{x}}$
\vspace{0.3in}

%\part $\displaystyle x^{\frac 1 2}-y^{\frac 1 3}=x$
%\vspace{2in}
\end{parts}

%\question Find the maxima, minima, and points of inflection of
%$f(x)=x^3+3x^2-9x+2$.  Sketch the graph.
%\vspace{3in}

\question Evaluate the following integrals.
\begin{parts}

\part $\displaystyle \int \sqrt{x} + x^{-\frac 3 2}\, dx$
\vspace{0.3in}

\part $\displaystyle \int \frac {\sin x \, dx}{(1+2\cos x)^2}$
\vspace{0.3in}


\part $\displaystyle \int \sec^2(x) \sec(\tan x)\;dx$
\vspace{0.3in}

\end{parts}










\end{questions}
\end{document}